% tla_bash.tex - 生成TLC终端输出的PDF图片（最终完美版）
% 核心：添加varwidth选项，支持多行内容，解决LR mode报错
\documentclass[varwidth]{standalone}
\usepackage{minted}  % 终端输出语法高亮
\usepackage{xcolor}  % 颜色配置
\usepackage[T1]{fontenc} % 字体编码兼容
\usepackage{lmodern} % 现代字体兼容

% 仅保留所有版本通用的核心配置，避免任何兼容问题
\setminted{
    fontsize=\large,       % 放大终端字体，便于论文中阅读
    linenos=false,         % 终端输出不显示行号
    breaklines=true,       % 自动换行，避免内容溢出
    bgcolor=gray!5,        % 浅灰背景，美观且不影响打印
    fontfamily=courier,    % 终端等宽字体，字符对齐
    encoding=utf8          % 支持UTF-8编码
}

\begin{document}
% 终端命令 + 输出内容（原样保留）
\begin{minted}{text}
>start SER
Tx_mulberry-server-2_1709561234890> read account-1
<nil>
Tx_mulberry-server-2_1709561234890> write account-1 100
OK
Tx_mulberry-server-2_1709561234890> commit
Committed
\end{minted}
\end{document}
